\textbf{Problem 3(c) answer}\GradescopeComment{This should be on page 7 of your PDF.}

\begin{enumerate}
  
\item[(i)] Define $C(n)$ to be the number of letters printed by {\sc PrintCs} given input $n\ge 1$.

  Then $C(1) = \BLUE{
    \REPLACEME 
  }$

  For $n>1$, $C(n) = \BLUE{
    \REPLACEME 
  }$
  
\item[(ii)] Within level $i$:

  The number of nodes is $\BLUE{
    \REPLACEME 
  }$
  
  The problem size is $\BLUE{
    \REPLACEME 
  }$

  The number of letters printed by each call is $\BLUE{
    \REPLACEME 
  }$

\item[(iii)] The number of levels is $\BLUE{
    \REPLACEME 
  }$
  
\item[(iv)] The total number of letters printed is 
  \[\BLUE{ C(n) ~=~ 
      %
      \REPLACEME
      %
      % \sum_{i=0}^{...}...
      %
    .}
  \]

\item[(v)] Using
  \BLUE{
    \REPLACEME
  },
  the big-$\Theta$ value of $C(n)$ is
  \[\BLUE{
      % 
      \REPLACEME
      % \Theta(...)
      % 
    .}\]

\end{enumerate}

\vfill

\paragraph{Problem \arabic{problem} collaboration and citations}\GradescopeComment{This should be on page 7 (bottom) of your PDF.}

I collaborated on this problem with the following people:
\begin{blue}
  \begin{itemize}
  \item \REPLACEME                    % write "none" if none
  \end{itemize}
\end{blue}

My answers use ideas from the following sources (other than the lecture notes and textbooks): 
\begin{blue}
  \begin{itemize}
  \item \REPLACEME                    % write "none" if none
  \end{itemize}
\end{blue}


%%% Local Variables:
%%% mode: latex
%%% TeX-master: "homework_2"
%%% End:
